\section{Conclusion et perspectives}
\label{sec:conclusion}

Ce projet a permis de mettre en \oe uvre l'essentiel de la cha\^ine d'un syst\`eme RAG appliqu\'e
aux concepts NLP/ML : constitution d'un corpus Wikipedia, pr\'eparation des donn\'ees, indexation
vectorielle avec ChromaDB, pipeline de r\'ecup\'eration et de g\'en\'eration via Groq, \'evaluation
quantitative sur 34 questions et visualisation des performances.

Les r\'esultats obtenus (precision@3 = 0.922, recall@3 = 0.985, fid\'elit\'e = 0.887, temps
moyen = 0.811~s) montrent que le syst\`eme r\'epond de mani\`ere pertinente et rapide sur
l'ensemble du corpus choisi. L'interface Streamlit a \'et\'e valid\'ee par des tests interactifs confirmant la qualit\'e
des r\'eponses en conditions d'utilisation r\'eelle. La documentation (\texttt{README.md}) et
le script de d\'emonstration (\texttt{DEMO.md}) accompagnent la soutenance du 15 juin 2026.

\subsection{Perspectives}

\begin{itemize}[noitemsep]
  \item Remplacer la m\'etrique de fid\'elit\'e lexicale par une approche LLM-as-judge.
  \item Ajouter un re-ranking (cross-encoder) pour am\'eliorer la precision sur les sujets
        recouvrants.
  \item \'Etendre le corpus (sources r\'ecentes, contenu multilingue).
  \item Permettre la citation explicite des sources dans le texte de la r\'eponse g\'en\'er\'ee.
\end{itemize}

Le code source, les scripts d'\'evaluation et ce rapport LaTeX sont disponibles sur le d\'ep\^ot :
\url{https://github.com/Muhammed-OTP/rag-nlp-m1}.
